\section{Sharding a simulation across multiple GPUs}
\label{architecture-multigpu}

This section is a design analysis of work Nimbus would need. It does not
describe an implemented distributed solver. The current backend configuration
selects one device index for a compute runtime. The
\href{../vulkan/include/compute/vulkan_device_context.hpp}{Vulkan context}
owns one selected physical device, logical device and compute queue; the
\href{../metal/src/compute/internal/metal_device_context.hpp}{Metal context}
owns one device session and command queue. There is no maintained MPI domain
decomposition or complete multi-device timestep owner in the ordinary runtime.
Multiple runtimes, nested domains, CUDA-assisted kernels and remote SSH
execution do not, individually or together, establish that capability.

\subsection{What is being divided?}

There are three substantially different designs. Independent simulations can
run on different devices without exchanging atmospheric state. Spatial
decomposition instead divides one atmosphere into owned subdomains, which
must exchange the state needed by neighboring numerical operators. Assigning
different physics packages to different GPUs is a third design: for example,
dynamics on one device and radiation on another. That arrangement must move
inputs and return tendencies at their declared cadence. It may duplicate much
of the state and adds a coupling dependency, so it should not be assumed to
solve a memory-capacity problem.

For an initial spatial implementation, horizontal tiles with complete vertical
columns are a reasonable candidate. They keep column radiation, boundary-layer
solves, vertical sedimentation and a column-local implicit solve on one owner.
This is a proposed starting restriction, not proof of support for every current
column formulation. It can produce poor aspect ratios or leave too few tiles
for a large machine. Splitting vertically later would require explicit
distributed column algorithms or gather/scatter phases, including their
memory and synchronization costs. Terrain and moving nests make ownership
more complex even when columns remain whole.

\begin{figure}[htbp]
\centering
\begin{tikzpicture}[
  box/.style={draw=NimbusBlue,fill=NimbusLight,rounded corners=2pt,
    text width=4.9cm,minimum height=17mm,align=center,font=\small},
  arrow/.style={-{Latex[length=2mm]},draw=NimbusBlue,thick}]
  \node[box] (a) at (0,0) {GPU A: owned horizontal tile\\complete vertical columns};
  \node[box] (b) at (7.2,0) {GPU B: owned horizontal tile\\complete vertical columns};
  \draw[arrow] ([yshift=3mm]a.east)--node[above,font=\scriptsize]{stage $s$ halo}([yshift=3mm]b.west);
  \draw[arrow] ([yshift=-3mm]b.west)--node[below,font=\scriptsize]{stage $s$ halo}([yshift=-3mm]a.east);
  \node[box] (stage) at (3.6,-2.4) {Validate the shared stage\\global timestep and owned budgets};
  \draw[arrow] (a.south)--(stage.north west);
  \draw[arrow] (b.south)--(stage.north east);
  \node[box,below=8mm of stage] (commit) {One accepted endpoint\\consistent output and checkpoint};
  \draw[arrow] (stage)--(commit);
\end{tikzpicture}
\caption{Proposed two-tile solver decomposition. Halo exchange can use a
validated direct path or explicit host staging. Arrows describe data
dependencies, not a mandatory global barrier after every local operator.
Each tile owns its physical cells once; copied halos carry no additional
atmospheric inventory. This is a design diagram, not a current Nimbus feature.}
\label{fig:architecture-multigpu}
\end{figure}


\subsection{Own physical cells and faces before assigning buffers}

Each tile needs global physical bounds, local storage bounds, halo widths,
neighbors and a stable mapping from global to local indices. A halo is a
readable copy of another tile's state, not additional atmospheric volume.
Inventory reductions must count owned cells once. On a C-grid, a velocity or
mass-flux face shared by two tiles also needs a canonical owner or an equivalent
reconciliation rule. Independently evaluating slightly different fluxes on
both sides can create an interface source even when each local kernel passes
its own conservation test.

For one mass flux $\Phi_f$ oriented from tile $A$ to tile $B$, conservative
exchange requires
\[
 \Delta M_A=-\Delta t\,\Phi_f,\qquad
 \Delta M_B=+\Delta t\,\Phi_f.
\]
The same face, time interval, stage weight and transported carrier must be used
on both sides. Water mixing ratios, number per dry mass and number per volume
cannot share an unspecified scalar convention. Geometry, metric coefficients,
held forcing and scheme-private fields require exchange wherever their
consumers need them. A fixed list of visible plot fields is not a sufficient
halo protocol.

Messages should identify the grid or hierarchy generation, physical interval,
integrator stage, field and location, source and destination ownership, extent
and precision. A completed transfer with the wrong stage is stale data. Halo
width must follow the largest actual dependency footprint of the selected
operator, including reconstruction and corner dependencies; it cannot be
guessed from the nearest-neighbor appearance of the final divergence.

Physical boundary conditions apply at the atmosphere's external boundary.
An internal partition face must not become a wall, sponge or independent open
boundary. Periodic neighbors must close the global topology. A processor
partition on a fixed grid is also different from a coarse/fine nest: the
latter needs metric-aware restriction, prolongation and flux reconciliation,
not merely a copy between equal-resolution halos.

\subsection{Exchange at the stage that consumes the data}

A portable first implementation could establish this sequence for each
operation with neighbor dependencies: complete the producing work, pack the
required owned boundary state, start its transfers, compute independent
interior cells, complete and unpack the receives, then compute the boundary
region. The new endpoint becomes visible only after all required regions
and validations complete. Packing, transfer and unpacking each retain buffer
ownership until their respective completion events.

This sequence must be derived from the actual split integrator. Acoustic
substeps may require pressure and face-wind exchange several times during one
slow step. Runge--Kutta stages consume their own stage state. Exchanging only
after a complete outer timestep is generally insufficient. A global CFL
decision must use the most restrictive participating tile; a local reduction
does not establish the safety of the global timestep. Iterative diagnostics
or solvers can add global reductions beyond those required by local stencils.

The current \href{../include/compute/dispatch_graph.hpp}{dispatch graph}
describes field hazards, but is not a distributed execution engine. It would
need region and device ownership, explicit exchange edges and completion
dependencies. Likewise, the
\href{../include/compute/compute_residency.hpp}{residency contract} would need
to distinguish an owned interior version from a halo version for a particular
stage. Marking an entire array ``current on the device'' cannot express that
one boundary slab has not arrived.

\subsection{Portable communication first, optional direct paths second}

The reference transport should work with explicit device-to-host transfer,
owned host send/receive buffers and host-to-device transfer. This gives a
testable route without assuming peer access. On a single machine the exchange
can be managed within one process or across explicitly owned processes.
Across machines, MPI is a plausible transport and collective layer. Its
\href{https://www.mpi-forum.org/docs/mpi-4.1/mpi41-report/node235.htm}{nonblocking
neighborhood operations} provide an interface for neighbor exchange, but
posting a nonblocking operation is not proof of overlap or completion. Progress,
thread support and the application's wait schedule must be measured and
validated. The transport does not know Nimbus's field units or stage semantics.

Direct device transfers are optional optimizations of the same exchange
contract. Vulkan device groups are enumerated compatibility groups; they do
not promise that arbitrary installed GPUs can be treated as one memory pool.
The \href{https://docs.vulkan.org/guide/latest/extensions/device_groups.html}{Khronos
device-group guide} and
\href{https://docs.vulkan.org/refpages/latest/refpages/source/vkGetDeviceGroupPeerMemoryFeatures.html}{peer-memory
query} describe the relevant capability checks. Separate logical devices
instead require a supported sharing/transfer mechanism with explicit lifetime
and synchronization rules. A queue barrier or fence on one independent device
does not by itself establish that another device can consume the bytes.

Metal similarly requires using the devices actually exposed by the system.
Apple documents
\href{https://developer.apple.com/documentation/metal/multi-gpu-systems}{multi-GPU
systems} and
\href{https://developer.apple.com/documentation/metal/transferring-data-between-connected-gpus}{transfers
between connected GPUs}. Availability depends on hardware and API support;
multiple packages or dies inside a product do not imply multiple independently
selectable Metal devices. Unified CPU/GPU memory on one device does not establish
peer access to another computer or accelerator.

``GPU-aware MPI'' also needs a specific memory API, transport and build. For
example, \href{https://docs.open-mpi.org/en/v5.0.8/tuning-apps/networking/cuda.html}{Open
MPI's CUDA support} describes supported CUDA-buffer paths. It does not certify
arbitrary Vulkan allocations or Metal buffers. Nimbus would need a validated
interoperability path or host staging for those owners. An unavailable direct
path may select an explicitly supported staged transport; it must not silently
change the numerical operation or violate a native-execution requirement.

\subsection{Why more memory does not imply linear speedup}

Let an illustrative cubic owned tile have side $n$, halo depth $h$, $F$
communicated scalar-equivalents and $b$ bytes per scalar. Its owned payload is
$bFn^3$. A complete surrounding halo occupies
\[
 B_{\mathrm{halo}}=bF\big[(n+2h)^3-n^3\big]
 =bF(6hn^2+12h^2n+8h^3).
\]
If only the six face slabs are communicated in an exchange, their leading
payload is $6bFhn^2$; corners, staged directional exchanges and non-scalar
locations change the exact count. This is a teaching model for an interior
tile, not the traffic of a particular Nimbus kernel. As a fixed global problem
is divided into more pieces, the surface-to-volume ratio grows. Full-column
horizontal tiles have a different geometric count but the same concern.

For message startup cost $\alpha$, effective transfer bandwidth $\beta$,
$m$ messages and payload $B$, a simple model is
\[
 T_{\mathrm{exchange}}\simeq m\alpha+B/\beta
       +T_{\mathrm{pack}}+T_{\mathrm{unpack}}.
\]
With perfect interior/exchange overlap, and with boundary work and a collective
following those phases, a stage still takes roughly
\[
 T_{\mathrm{stage}}\gtrsim
 \max(T_{\mathrm{interior}},T_{\mathrm{exchange}})
 +T_{\mathrm{boundary}}+T_{\mathrm{collective}}.
\]
Real overlap can be worse: packing can contend for memory bandwidth, transfer
engines may share links, and small tiles can underfill a GPU. Frequent acoustic
exchanges can become latency-bound. A second board behind a narrow or shared
link may increase capacity while adding little useful throughput.

Every device must fit its own owned state, halos, stages, rollback and
communication buffers. The sum of nominal VRAM capacities is not an admission
test. Nor should every rank retain a full global host mirror: that would
preserve a memory bottleneck even after the device fields were partitioned.
NUMA placement should associate host staging buffers and communication threads
with the relevant device and network path. On heterogeneous GPUs, equal cell
counts need not give equal stage times. Cloud-dependent work and nested
placement can also make the slowest tile change during an experiment.
Dynamic balancing would require a state migration transaction and new
generations, not just moving an array pointer.

\subsection{Global diagnostics, restart and failure must agree}

Budget reductions should accumulate owned physical inventories with declared
precision and ordering. A reduction tree may change rounding when the number
of ranks changes. Fixed-decomposition repeatability, decomposition-independent
bitwise equality and bounded scientific agreement are distinct claims. They
need distinct tests; ordinary parallel summation does not promise all three.
Random perturbations should have a defined mapping to physical identities and
time if results are intended to be independent of rank placement.

A checkpoint needs a globally consistent accepted endpoint: all owned primary
and private fields, geometry, held-source clocks, histories, configuration and
publication identity. Repartitioned restart needs a mapping from saved physical
ownership to the new layout; it cannot depend on finding the same GPU index.
Partitioned output should publish a manifest only after all required pieces
are durable, and readers should reconstruct the logical field without treating
rank order as a coordinate system. Gathering the whole atmosphere on one rank
for every frame can become a new memory and I/O bottleneck. Existing NPY,
netCDF and ZFP support does not by itself supply this distributed publication
protocol, and this design does not authorize removal of NPY.

If one tile rejects a timestep, neighboring tiles cannot publish a mixture of
old and new endpoints. A first design should stop the distributed job on an
unrecoverable device or communication failure and resume from a validated
global checkpoint. Transparent continuation after losing a rank is additional
engineering, not a consequence of using MPI. Buffers involved in uncertain
device submissions must remain owned until completion is known or the
corresponding execution context has been safely retired.

\subsection{A sequence that can be implemented and reviewed}

Begin with the distributed state and exchange contract, including a one-tile
path that reproduces the existing admitted experiment. A two-tile CPU or
host-staged reference can then expose ownership errors independently of peer
transport. Manufactured fields should cross each partition orientation and
corner; interface inventories must cancel, and halos must never enter physical
budgets. Tests should mutate stage identities, interrupt transfers and reject
partial publication.

The next bounded target is a two-device, fixed-grid configuration with stated
physics and boundary conditions, selected through ordinary public configuration.
It needs whole-step one-device comparisons, native dispatch evidence, restart
on the same decomposition, rejected-input atomicity and memory accounting.
Changing the decomposition on restart, nonuniform hardware, terrain, private
microphysics state, held radiation, nesting and moving domains each extend the
contract and require their own completed coupling. They cannot be called
supported merely because the flat dry exchange works.

Only after correctness should direct peer transfers, overlap, larger device
counts and multi-node transports be assessed with strong- and weak-scaling
experiments. Strong scaling holds the physical numerical problem fixed; weak
scaling holds approximately the work per participant fixed and generally grows
the global problem. Report both completion time and resource cost, including
output and restart. Source code or a two-GPU speedup alone is not Ready, and
Ready remains separate from campaign validation and independent scientific
review. This staged design keeps the ordinary local workflow usable while
making distributed execution an explicit, testable extension.
